\documentclass[10pt]{article}
\usepackage[T1]{fontenc}
\usepackage[utf8]{inputenc}
\usepackage[french]{babel}
\usepackage{newpxtext,newpxmath}
\usepackage[a4paper,top=1.5cm,bottom=1.35cm,left=1.6cm,right=1.6cm]{geometry}
\usepackage{microtype}
\usepackage{xcolor}
\usepackage{titlesec}
\usepackage[most]{tcolorbox}
\usepackage{enumitem}
\usepackage{array}
\usepackage{booktabs}
\usepackage{fancyhdr}

\definecolor{navy}{HTML}{1F3A5F}
\definecolor{grey}{HTML}{54606E}
\definecolor{green}{HTML}{2C6E49}
\definecolor{amber}{HTML}{9C6A16}
\definecolor{red}{HTML}{A6322E}
\definecolor{plum}{HTML}{7A4EA8}
\definecolor{boxgray}{HTML}{F3F5F8}
\definecolor{boxred}{HTML}{FCF2F1}
\definecolor{boxgreen}{HTML}{F1F7F3}
\definecolor{boxplum}{HTML}{F7F3FB}
\definecolor{boxcode}{HTML}{F7F8FA}
\definecolor{boxedge}{HTML}{D9E0EA}

\setlength{\parindent}{0pt}
\setlength{\parskip}{2.6pt}
\linespread{1.0}

\pagestyle{fancy}\fancyhf{}
\lhead{\small\color{grey}Scella}
\rhead{\small\color{grey}Le groupe le plus révolutionnaire -- 11 août 2026}
\cfoot{\small\color{grey}\thepage}
\renewcommand{\headrulewidth}{0.3pt}

\titleformat{\section}{\color{navy}\large\bfseries}{}{0pt}{}
\titlespacing{\section}{0pt}{11pt}{3pt}
\titleformat{\subsection}{\color{navy}\bfseries}{}{0pt}{}
\titlespacing{\subsection}{0pt}{7pt}{1.5pt}

\newtcolorbox{these}{colback=boxgray,colframe=boxedge,boxrule=0.5pt,arc=2.5pt,
  left=8pt,right=8pt,top=4.5pt,bottom=4.5pt}
\newtcolorbox{alerte}{colback=boxred,colframe=red!55,boxrule=0.8pt,arc=2.5pt,
  left=8pt,right=8pt,top=4.5pt,bottom=4.5pt}
\newtcolorbox{acquis}{colback=boxgreen,colframe=green!45,boxrule=0.8pt,arc=2.5pt,
  left=8pt,right=8pt,top=4.5pt,bottom=4.5pt}
\newtcolorbox{piste}{colback=boxplum,colframe=plum!45,boxrule=0.8pt,arc=2.5pt,
  left=8pt,right=8pt,top=4.5pt,bottom=4.5pt}
\newtcolorbox{extrait}{colback=boxcode,colframe=boxedge,boxrule=0.5pt,arc=1.5pt,
  left=7pt,right=7pt,top=3.5pt,bottom=3.5pt}

% \textcolor plutôt que {\color …} : en tête de cellule p{}, \color insère un whatsit
% qui ouvre une ligne vide et décale verticalement toute la colonne.
\newcommand{\best}[1]{\textcolor{green}{\textbf{#1}}}
\newcommand{\pire}[1]{\textcolor{red}{\textbf{#1}}}
\newcommand{\sz}[1]{\textcolor{grey}{\small\textsf{#1}}}
\newcommand{\code}[1]{\texttt{\small #1}}
\newcommand{\x}{$\times$}
\newcommand{\loc}[1]{{\footnotesize\color{grey}\texttt{#1}}}

\begin{document}

\begin{center}
{\Huge\bfseries\color{navy} Scella}\\[4pt]
{\large\color{grey}\itshape Le groupe le plus révolutionnaire}\\[4pt]
{\small\color{grey} 11 août 2026 \; · \; axe orthogonal au classement produit \; · \; six levées d'impossibilité, localisées}
\end{center}
\vspace{2pt}

\begin{these}
\textbf{Le résultat central : les crates ne changent presque pas, la thèse change entièrement.} Le
dépôt contient \textbf{six levées d'impossibilité}, et j'ai vérifié où chacune vit :
\textbf{cinq sur six sont déjà co-localisées} dans les neuf crates du produit désigné la veille. La
sixième n'en coûte qu'\textbf{un seul de plus}. Mais l'analyse précédente, qui optimisait la
\emph{cohérence produit}, avait retenu une promesse plus faible que son contenu -- « l'essaim \textbf{attestable} », c'est-à-dire
\emph{le système peut prouver qu'il a raison}. Sur l'axe révolutionnaire, la promesse juste est son
\textbf{inverse épistémique} : \emph{le système peut être \textbf{pris en défaut}} -- par lui-même au
démarrage, par n'importe qui ensuite sans clé ni dépendance, et il dit lui-même comment le défier. La
première catégorie est saturée ; la seconde est presque vide.
\end{these}

\section{1 · Ce que « révolutionnaire » exige, défini pour être vérifiable}

Le mot ne vaut rien s'il reste une appréciation, et il serait facile de s'en servir pour re-couronner
la conclusion de la veille. Je l'opérationnalise donc, autour d'un test \textbf{binaire} :

\begin{extrait}
\textbf{Quelle phrase devient vraie pour la première fois ?} \\[2pt]
Si la phrase est « X est 30 \% plus rapide, plus propre, moins cher », c'est une \emph{amélioration}. \\
Si la phrase est « X peut désormais être fait, \textbf{par quelqu'un qui ne le pouvait pas} », c'est une
\textbf{levée}.
\end{extrait}

Cinq critères en découlent. L'\textbf{impossibilité levée} -- binaire, jamais graduelle. Le
\textbf{transfert de pouvoir} : une révolution déplace un pouvoir, pas une métrique, et la question
utile est \emph{qui peut désormais contrôler ce qu'il devait croire ?} L'\textbf{irréversibilité} :
une fois la chose vue, l'attente par défaut est-elle changée ? L'\textbf{absence d'équivalent} -- non
pas « rare », ce qui était la lentille précédente, mais « personne ne sait le faire ». Et la
\textbf{portée} : cela dépasse-t-il le cas d'usage du dépôt ?

\begin{alerte}
\textbf{Cet axe déclasse explicitement ceux qui portaient le verdict de la veille.} La
\textbf{maturité} et la \textbf{démontrabilité} avaient désigné l'essaim attestable parce qu'il tourne
en une commande aujourd'hui. Ici elles ne comptent presque pas : une levée d'impossibilité reste une
levée même si elle se montre mal. C'est précisément pour cela que le résultat pouvait différer -- et
qu'il fallait refaire le travail plutôt que défendre la conclusion.

Je maintiens en revanche la contrainte constante de la série : \textbf{aucune cible extérieure n'est
déclarée, et je n'en invente pas}. Les cinq critères se jugent sans destinataire.
\end{alerte}

\section{2 · Les six levées, localisées dans le code}

\renewcommand{\arraystretch}{1.22}
\begin{center}\small
\begin{tabular}{@{}c >{\raggedright\arraybackslash}p{3.15cm} >{\raggedright\arraybackslash}p{6.55cm} >{\raggedright\arraybackslash}p{5.0cm}@{}}
\toprule
& \textbf{La levée} & \textbf{La phrase qui devient vraie} & \textbf{Où} \\
\midrule
\textbf{A} & Vérifieur zéro-dépendance & L'auditeur lit l'\textbf{intégralité} du vérifieur avant de lui accorder sa confiance & \code{verify-meta-mini} \sz{arbre : 1 ligne} \\
\textbf{B} & Ancre $\rightarrow$ troncature sans clé & Un tiers ne détenant \textbf{aucune} clé détecte qu'on lui a montré la moitié qui arrange & \code{provenance-dag/anchor.rs} \newline \code{controller/anchor.rs} \\
\textbf{C} & Liaison v4 nominative & Renommer un acteur \textbf{casse la preuve} & \code{efe-meta-verify/src/lib.rs} \\
\textbf{D} & Canari + porte H1 & Le système \textbf{mesure} sa garantie, \textbf{se réfute}, et \textbf{refuse de démarrer} & \code{participants/lipschitz.rs} \newline \code{participant/local.rs} \\
\textbf{E} & Triangle croisé & Deux implémentations doivent s'accorder $\rightarrow$ le format devient une \textbf{spécification} & \code{verify-meta-mini/tests/} \newline \code{cross\_check.rs} \\
\textbf{F} & Frontière PROUVÉ / CONFIANCE & L'artefact dit lui-même \textbf{comment le défier} & \code{forensic-export/src/} \newline \code{readme.rs} \\
\bottomrule
\end{tabular}
\end{center}

Deux d'entre elles doivent être lues avant d'être notées, parce que leur formulation abrégée ne rend
pas ce qu'elles font.

\subsection{2.1 · D -- le système qui réfute sa propre documentation}

Le document d'architecture de la Couche 3 porte, \emph{dans son propre texte}, la réfutation de son
propre certificat : \emph{« le chiffre 0.17 est RÉFUTÉ (mesure : $\hat{L} = 0{,}366$) »}, l'ancien
nombre barré partout où il apparaissait, et -- c'est la partie qui compte -- la séparation explicite
de ce qui survit et de ce qui tombe : \emph{« Φ reste largement contractant ($0{,}366 \ll 1$, H1
tient) »} d'un côté, \emph{« la forme fermée est réfutée comme certificat »} de l'autre.

Un système, normalement, \textbf{affirme} ses propriétés de sûreté dans sa documentation, et la
documentation fait foi jusqu'à preuve du contraire. Celui-ci les \textbf{mesure}, se contredit, retient
la mesure défavorable, et pose une porte de démarrage qui refuse de servir si la mesure sort du
domaine. Le rapport de force entre la fiche technique et le binaire est inversé.

\subsection{2.2 · F -- l'artefact de conformité qui argumente contre lui-même}

\code{forensic-export} \textbf{génère, à l'intérieur du bundle livré}, un paragraphe destiné à
l'auditeur. Quand l'ancrage est purement interne, il écrit :

\begin{extrait}
\sz{« … une \textbf{déclaration interne}, \textbf{sans force probante externe} par construction. Elle
\textbf{n'établit PAS} que le record n'a pas été fabriqué a posteriori (le déclarant et le détenteur de
la clé peuvent être la même partie). Pour une preuve \textbf{opposable au détenteur}, exigez un ancrage
à \textbf{reçu externe vérifiable} (RFC 3161 / log de transparence / dépôt régulateur). »}
\end{extrait}

Et quand il n'y a aucun ancrage, il le dit aussi -- avec une section intitulée littéralement
\mbox{\textbf{« Ce que ça NE prouve PAS »}}.

C'est un artefact de conformité qui \textbf{instruit son lecteur sur la manière de le défaire}, et qui
\textbf{nomme le remède à exiger} du fournisseur -- c'est-à-dire de lui-même. Toutes les incitations
du domaine vont dans l'autre sens : un artefact de conformité est habituellement écrit pour clore la
discussion, pas pour l'armer.

\section{3 · Le classement sur les cinq critères}

Rangs de 1 à 6 par critère ; \emph{Éc.} est l'amplitude entre le meilleur et le pire rang.

\renewcommand{\arraystretch}{1.16}
\begin{center}\small
\begin{tabular}{@{}r >{\raggedright\arraybackslash}p{7.0cm} ccccc c c@{}}
\toprule
& \textbf{La levée} & \sz{Imposs.} & \sz{Pouvoir} & \sz{Irrév.} & \sz{Sans éq.} & \sz{Portée} & \textbf{Moy} & \sz{Éc.} \\
\midrule
\textbf{1=} & \textbf{B · Ancre sans clé} & \best{1} & \best{1} & 4 & 4 & 2 & \textbf{2,40} & 3 \\
\textbf{1=} & \textbf{F · L'artefact qui dit comment le défier} & 3 & 2 & 2 & 2 & 3 & \textbf{2,40} & \best{1} \\
\textbf{3=} & \textbf{A · Vérifieur zéro-dépendance} & 4 & 3 & 3 & 3 & \best{1} & \textbf{2,80} & 3 \\
\textbf{3=} & \textbf{D · Auto-réfutation + refus de démarrer} & 2 & 4 & \best{1} & \best{1} & \pire{6} & \textbf{2,80} & \pire{5} \\
5 & C · Liaison v4 nominative & 5 & 5 & 5 & \pire{6} & 4 & \textbf{5,00} & 2 \\
6 & E · Triangle croisé & \pire{6} & \pire{6} & \pire{6} & 5 & 5 & \textbf{5,60} & \best{1} \\
\bottomrule
\end{tabular}
\end{center}

Trois lectures comptent davantage que les rangs eux-mêmes.

\textbf{B est premier là où c'est le plus difficile} -- impossibilité \emph{et} transfert de pouvoir.
C'est la seule levée du dépôt qui fonctionne \textbf{contre le détenteur de la clé lui-même}. Une
chaîne HMAC ne peut pas, \emph{par construction}, détecter sa propre troncature : le préfixe d'une
chaîne valide est une chaîne valide. Tout ce que la cryptographie du reçu peut faire s'arrête là.
L'ancre est le seul mécanisme qui retire un pouvoir à celui qui détient le secret, et le contrôle qui
l'exerce ne demande aucune clé -- c'est une comparaison de chaînes de caractères.

\textbf{F arrive premier ex \ae quo avec un écart de 1} : bon partout, faible nulle part. C'est le
profil d'une levée \emph{structurelle} plutôt que spectaculaire, et c'est la seule des six à se
trouver \textbf{hors} du produit retenu la veille.

\textbf{D porte un écart de 5} -- premier en irréversibilité et en absence d'équivalent,
\textbf{dernier en portée}. La chose la plus inimitable du dépôt ne s'applique qu'aux systèmes
numériques contractants. C'est exactement le profil haute-variance déjà diagnostiqué au niveau des
ensembles, retrouvé ici sur un axe entièrement différent : une convergence entre deux méthodes
indépendantes, qui rend le diagnostic plus crédible que chacune prise seule.

\sz{Et E tombe dernier, ce qui est cohérent : le triangle croisé rend la correction plus \emph{probable},
il ne rend rien \emph{possible}. C'était une force au classement produit ; c'en est une faible ici. Les
deux axes ne mesurent pas la même chose, et il faut l'accepter plutôt que le lisser.}

\section{4 · Le groupe -- et le fait qu'il était déjà là}

La vérification des emplacements donne le résultat le plus utile du document.
\textbf{Cinq levées sur six sont déjà co-localisées} dans les neuf crates du produit désigné la veille.
La sixième, F, coûte \code{forensic-export}, dont la fermeture mesurée est de deux crates --- mais
le second est \code{provenance-dag}, \textbf{déjà présent} dans le produit. Le coût marginal réel est
donc d'\textbf{un seul crate}.

\begin{acquis}
\textbf{Le groupe le plus révolutionnaire} $=$ A $+$ B $+$ C $+$ D $+$ E $+$ F, soit
\textbf{dix crates} sur soixante-six. Les crates ne changent quasiment pas.
\textbf{C'est la promesse qui change} -- et c'est elle qui sépare « bien empaqueté » de
« révolutionnaire ».
\end{acquis}

\subsection{4.1 · Le renversement de la promesse}

C'est le résultat central, et il corrige l'analyse précédente sans l'annuler.

\renewcommand{\arraystretch}{1.2}
\begin{center}\small
\begin{tabular}{@{}l >{\raggedright\arraybackslash}p{5.0cm} >{\raggedright\arraybackslash}p{8.0cm}@{}}
\toprule
& \textbf{La promesse} & \textbf{Ce qu'elle revendique} \\
\midrule
Avant & « L'essaim \textbf{attestable} » & \emph{le système peut prouver qu'il a raison} \\
\textbf{Après} & « Le système \textbf{réfutable} » & \emph{le système peut être \textbf{pris en défaut}} \\
\bottomrule
\end{tabular}
\end{center}

La première est une assertion de fiabilité : la catégorie la plus encombrée qui soit, et celle où
chaque affirmation ressemble à celle du voisin. La seconde est son \textbf{inverse épistémique}, et
elle est presque vide.

\begin{piste}
\textbf{« Un système qui peut être pris en défaut : par lui-même au démarrage, par n'importe qui
ensuite -- sans clé et sans dépendance -- et qui dit lui-même comment le défier. »}
\end{piste}

Les quatre levées de tête s'y rangent sans forcer, et c'est le signe qu'il s'agit bien d'une thèse et
non d'un slogan : \textbf{D} le prend en défaut au démarrage, \textbf{B} permet à un tiers sans clé de
le prendre en défaut ensuite, \textbf{A} permet d'auditer l'auditeur lui-même, et \textbf{F} apprend au
lecteur comment s'y prendre. L'analyse de la veille avait choisi les bons crates pour une raison plus
faible que ce qu'ils contiennent.

\section{5 · Le second groupe, et il faut le dire honnêtement}

Si « révolutionnaire » autorise ce qui n'est pas encore construit, la réponse change.

\textbf{L'interstice gouverné} -- Fil B (servi) $+$ Sandwich C.1 $+$ essaim médiateur.

\begin{piste}
\textbf{« L'espace entre deux modèles frontière n'est plus un tuyau opaque : c'est un lieu où des
agents déclarés opèrent, chacun scellé et attribuable. »}
\end{piste}

C'est la levée \textbf{conceptuellement la plus grande} du dépôt, et pour une raison de nature : elle
ne donne pas une propriété nouvelle à un objet existant, elle \textbf{crée un objet qui n'existait
pas} -- l'interstice comme lieu gouverné. Aucune des six levées précédentes ne fait cela ; toutes
améliorent le statut épistémique d'une chose déjà là.

Ce qui l'empêche de l'emporter est factuel et se rejoue en une commande : \code{HopKind},
\code{SwarmMediator} et \code{EdgePlan} comptent \textbf{zéro occurrence}. Le tiers servi -- Fil B, dont
le \code{slm\_mediated} vit dans \code{routes/llm.rs} -- est réel et complet ; les deux autres tiers
sont des documents d'architecture. Le substrat commun, lui, existe déjà : les deux binaires partagent
\code{provenance-dag} et le même \code{AnchorIndex}.

\textbf{Ma position, sans arrondi.} Le groupe réfutable est le plus révolutionnaire \textbf{parmi ce
qui existe} ; l'interstice gouverné l'est \textbf{en potentiel}, et son terrain est préparé. Je ne les
mets pas au même rang, parce qu'une révolution non écrite n'en est pas une -- mais je ne dissimule pas
que le second a le plafond le plus haut.

\section{6 · Ce que j'écarte, et pourquoi}

\textbf{SOMNIA-MLP} aurait le plafond le plus élevé de tous -- auto-adaptation embarquée certifiable,
Invariant 1 prouvé sous Frama-C. Je l'écarte pour des raisons déjà consignées et qui n'ont pas bougé :
il est \textbf{hors dépôt}, le matériel n'a \textbf{jamais été flashé}, le gold n'est pas résolu, et
empiler deux dépendances-gold non validées est exactement ce qu'il ne faut pas faire.

\textbf{La Colonne} reste écartée pour la raison du rapport précédent : sa phrase décrit une
\emph{architecture}, pas une levée -- un lecteur en ressort en sachant comment c'est bâti, pas ce qui
devient possible.

\section{7 · Le jugement}

Les crates bougent à peine : \textbf{dix au lieu de neuf}, l'unique ajout étant \code{forensic-export}. Ce qui
bouge est la \textbf{thèse}, et sur l'axe demandé c'est elle qui décide.

Le dépôt contient six levées d'impossibilité. Quatre sont de premier rang, et elles convergent toutes
vers la même propriété inhabituelle : \textbf{ce système est construit pour être attrapé}. C'est cela
qui n'a pas d'équivalent -- pas le consensus, pas le scellement, pas même le reçu pris isolément.

\textbf{Trois actions}, inchangées quant au code mais réordonnées quant au sens :
\textbf{absorber \code{forensic-export}} dans le produit -- un seul crate de plus, la sixième levée, et de loin
la moins chère ; \textbf{extraire la recette Mode-1 hors d'\code{efe-server}} -- mécanique, en
coordination avec la session concurrente ; \textbf{écrire le profil} qui allume les interrupteurs.

\begin{these}
\small\textbf{Rappel de cadrage, inchangé dans toute la série.} Le reçu atteste l'\textbf{intégrité}
d'une suite de pas, jamais la \textbf{correction} d'une décision. Il est \emph{tamper-evident pour un
détenteur de clé}, pas \emph{trustless pour un tiers} -- et c'est précisément pourquoi l'ancre (B) et
la frontière PROUVÉ/CONFIANCE (F) comptent autant : elles disent ce que le reçu seul ne peut pas
faire. Les participants du mode réel exécutent une fusion latente Mode-1 numérique : ce ne sont pas des
« vrais SLM ».
\end{these}

\section{Annexe · Vérifier chaque affirmation soi-même}

Les \emph{notes} des cinq critères sont un jugement, structuré pour être discutable et non mesurable.
Les \textbf{emplacements} et les \textbf{zéros}, eux, se rejouent.

\renewcommand{\arraystretch}{1.22}
\begin{center}\small
\begin{tabular}{@{}>{\raggedright\arraybackslash}p{5.3cm} >{\raggedright\arraybackslash}p{10.8cm}@{}}
\toprule
\textbf{Affirmation} & \textbf{Ce qui la rejoue} \\
\midrule
A · le vérifieur est sans dépendance &
  \code{cargo tree -p verify-meta-mini -e normal -{}-prefix none} \sz{(1 ligne)} \\
B · l'ancre sans clé &
  \code{grep -rl 'meta\_anchor\_of\textbackslash|META\_ANCHOR\_PREFIX' crates -{}-include='*.rs'} \\
C · la liaison nominative &
  \code{grep -rl 'struct RunBinding' crates -{}-include='*.rs'} \\
D · le canari et la porte &
  \code{crates/efe-meta-swarm/src/participants/lipschitz.rs} \newline
  \code{grep -rl 'fn h1\_boot\_gate' crates -{}-include='*.rs'} \\
D · l'auto-réfutation, écrite &
  \code{grep -n 'RÉFUT' docs/architecture/COUCHE3-M1-*.md} \\
E · le triangle croisé &
  \code{crates/verify-meta-mini/tests/cross\_check.rs} \\
F · l'artefact qui dit ses limites &
  \code{sed -n '28,56p' crates/forensic-export/src/readme.rs} \\
F · son coût de packaging &
  \code{cargo tree -p forensic-export -e normal -{}-prefix none} \newline \sz{2 crates, dont \code{provenance-dag} déjà dans le produit $\Rightarrow$ coût marginal $= 1$} \\
Le second groupe est à zéro &
  \code{grep -rc 'HopKind\textbackslash|SwarmMediator\textbackslash|EdgePlan' crates -{}-include='*.rs'} \\
Fil B, lui, est servi &
  \code{grep -rl 'slm\_mediated' crates -{}-include='*.rs'} \\
\bottomrule
\end{tabular}
\end{center}

Un avertissement pour finir, et il porte sur la nature même de ce document. Les six levées sont des
\textbf{faits} localisés ; leur \textbf{hiérarchie} est un jugement. Un désaccord sur « B avant F » est
légitime et ne se tranche par aucune commande -- mais un désaccord sur « \code{HopKind} : zéro
occurrence » s'ouvre en une seconde, et c'est cette asymétrie qui rend le classement discutable plutôt
qu'arbitraire.

\end{document}
